% Appendix C: Commitment Notes & Design Decisions (~12 pages)

\chapter{Commitment Notes \& Design Decisions}
\label{app:commitments}

이 부록은 SCC 이론의 핵심적 설계 결정들을 정당화와 함께 기록한다. 각 결정은 이론의 정체성을 규정하며, 일상적 수정의 대상이 아니다. 왜 이러한 결정이 내려졌는지를 이해하는 것은 이론을 올바르게 확장하는 데 필수적이다.


\section{존재론적 우선성: Soft Field가 원시이다}
\label{app:cn-ontological}

\begin{center}
\fbox{\parbox{0.85\textwidth}{\centering
\textbf{CN1. 존재론적 우선성}\\[4pt]
부드러운 응집 장 $u_t : X_t \to [0,1]$은 이론의 원초적 존재자이다.\\
이산적 객체는 항상 파생적이다. 이 우선성을 결코 역전시키지 않는다.
}}
\end{center}

\textbf{정당화.} 전-객체적 응집은 본질적으로 연속적이다. 형성이 인식 가능한 객체로 안정화되기 이전에, 응집은 정도를 허용한다: 어떤 영역은 강하게, 어떤 영역은 약하게 참여한다. 이진 집합 $A_t \subseteq X_t$로 시작하는 것은 이 연속적 전-객체 단계를 즉각 왜곡한다.

\textbf{결과.} 이산 부분집합, 클래스 라벨, 인스턴스 ID는 모두 응집 장에서 파생된다. 임계값 $\theta$를 적용하여 $A_t = \{x : u_t(x) \geq \theta\}$로 이산 객체를 복원할 수 있지만, 이 복원은 이차적 연산이지 기반으로의 회귀가 아니다.

\textbf{위반하면.} 이론을 이산 분할, 클러스터링, 추적으로 환원하면 전-객체적 상태---형성되고 있지만 아직 완전한 객체는 아닌 상태---를 기술하는 능력을 상실한다.


\section{층 분리: 존재론, 공리, 실현, 구현}
\label{app:cn-layers}

\begin{center}
\fbox{\parbox{0.85\textwidth}{\centering
\textbf{CN2. 층 분리}\\[4pt]
존재론, 공리, 연산자 실현, 구현의 네 층을 엄격히 구분한다.\\
어떤 층도 다른 층과 혼동하거나 붕괴시키지 않는다.
}}
\end{center}

네 층의 구조:

\begin{description}
    \item[존재론 (Ontology)] 이론이 원초적 존재로 상정하는 것. 응집 장 $u_t$, 닫힘 $\mathrm{Cl}_t$, 인접 $\mathbf{N}_t$, 구별 $\mathbf{D}_t$, 전달 $\mathbf{M}_{t \to s}$.

    \item[공리 (Axioms)] 연산자들이 만족해야 하는 보편적 성질. A1'(조건부 확장성), A2(단조성), A3(축소), B1--B3(인접), D-Ax1--3(구별), E1--E4(전달). 이들은 어떤 특정 실현에도 의존하지 않는다.

    \item[실현 (Realization)] 공리를 만족하는 구체적 함수 형태. 시그모이드 닫힘(Spec \S9.2), 시그모이드 구별(\S9.3), 분석해 공동-소속(\S9.4), Sinkhorn 전달(\S9.5). 이들은 현재 최선의 후보이지, 영구적 정의가 아니다.

    \item[구현 (Implementation)] \texttt{scc/} 패키지의 Python 코드. 이산화, 수치 정밀도, 최적화 알고리즘 등의 선택을 포함한다.
\end{description}

\textbf{정당화.} 층 분리는 이론의 견고성을 보장한다. 시그모이드 닫힘이 더 나은 실현으로 대체되더라도, 공리와 존재론은 변하지 않으며, 공리 위에 세워진 정리들도 유효하다.


\section{비멱등 닫힘: 의도적 설계}
\label{app:cn-nonidempotence}

\begin{center}
\fbox{\parbox{0.85\textwidth}{\centering
\textbf{CN3. 멱등성은 의도적으로 생략되었다}\\[4pt]
닫힘은 안정화 경향(A3: 축소)을 가지지, 원초적 멱등성을 가지지 않는다.\\
$\mathrm{Cl}(\mathrm{Cl}(u)) \neq \mathrm{Cl}(u)$는 이론의 기본 약속이다.
}}
\end{center}

\textbf{정당화.} 표준 위상학의 닫힘 연산은 멱등적이다: $\mathrm{Cl}(\mathrm{Cl}(A)) = \mathrm{Cl}(A)$. 그러나 전-객체 수준에서 닫힘은 이미 완료된 연산이 아니라 안정화 \emph{경향}이다. 응집 장은 ``거의'' 자기-지지적일 수 있지만 완벽하게 닫혀 있지는 않다.

\textbf{수학적 이점.}
\begin{enumerate}
    \item A3(축소)은 고정점의 유일성과 반복 수렴을 보장한다 ($a_{\mathrm{cl}} < 4$).
    \item Bind 진단 ($1 - \|u - \mathrm{Cl}(u)\|_2/\sqrt{n}$)은 자기-지지의 \emph{정도}를 연속적으로 측정한다. 멱등적 닫힘에서는 Bind가 자명하게 1이 되어 진단 가치를 상실한다.
    \item 닫힘 에너지 $\mathcal{E}_{\mathrm{cl}} = \|u - \mathrm{Cl}(u)\|^2$는 멱등적 닫힘에서 닫힘의 고정점에서만 0이 되어, 형성과 비형성을 이분법적으로 분류한다. 비멱등 닫힘에서는 연속적 잔차가 형성의 질을 정량화한다.
\end{enumerate}


\section{공리 vs.\ 실현: 보편성과 잠정성}
\label{app:cn-axiom-realization}

\begin{center}
\fbox{\parbox{0.85\textwidth}{\centering
\textbf{CN4. 공리는 보편적이고, 실현은 잠정적이다}\\[4pt]
A1'--E4는 보편적 공리이다. 시그모이드, 분석해, Sinkhorn은 현재 선택이다.\\
다른 실현도 가능하며, 공리를 만족하는 한 이론적으로 동등하다.
}}
\end{center}

\textbf{구체적 예.}
\begin{itemize}
    \item \textbf{닫힘}: 현재 시그모이드 $\sigma(a_{\mathrm{cl}}(\cdots))$. 대안: ReLU 기반, 다항식, 신경망 등. 공리 A1'--A3을 만족하면 모든 정리 유효.
    \item \textbf{구별}: 현재 시그모이드 비교 $\sigma(a_D(Pu - \lambda_D P(1{-}u)) - \tau_D)$. 대안: 정보-이론적 발산, 커널 기반 등. D-Ax1--3 만족 필요.
    \item \textbf{공동-소속}: 현재 분석해 $(I - \alpha W_{\mathrm{sym}})^{-1}$. 파생 진단이므로 실현 선택이 정리에 영향 없음.
    \item \textbf{전달}: 현재 Sinkhorn 부분 OT. 대안: 비정규화 OT, 확산 기반, 신경 ODE. E1--E4 만족 필요.
\end{itemize}


\section{$\mathbf{T}_t$와 $\mathbf{C}_t$의 강등}
\label{app:cn-demotion}

\begin{center}
\fbox{\parbox{0.85\textwidth}{\centering
\textbf{CN5. 사용되지 않는 것은 형식 우주에서 제거한다}\\[4pt]
$\mathbf{T}_t$ (v2.0 제거), $\mathbf{C}_t$ (v2.1 강등): 형식 우주에서\\
술어나 에너지에 기여하지 않는 구성요소는 파생 진단으로 강등한다.
}}
\end{center}

\subsection{전이 연산자 $\mathbf{T}_t$의 제거 (v2.0)}

$\mathbf{T}_t$는 경계 형태학을 특성화하기 위해 도입되었으나, 5차 반복 개발 동안:
\begin{itemize}
    \item 구체적 실현: 0개
    \item 증명된 정리: 0개
    \item 에너지 항: 0개
    \item 술어 역할: 0개
\end{itemize}

$\mathbf{T}_t$의 개념적 역할---경계 형태학 특성화---은 기울기 지표 $g_t$, 경계 대역 $\mathrm{Bd}_t$, 형태학적 질 측도 $\mathcal{Q}_{\mathrm{morph}}$에 의해 충분히 대체된다.

\subsection{공동-소속 연산자 $\mathbf{C}_t$의 강등 (v2.1)}

$\mathbf{C}_t$는 원래 분리 술어 $\mathsf{Sep}$에 진입했다. 그러나 $\mathsf{Sep}$가 $u$-가중 형태로 정정된 후:
\begin{itemize}
    \item 술어 역할: 0개 (Sep에서 제거됨)
    \item 에너지 항: 설계상 0개 (에너지 해석성 유지)
\end{itemize}

$\mathbf{C}_t$의 $\mathbf{C}_t$-가중 Sep 공식은 형성의 질에 무관하게 항상 $\approx 0.5$를 주었다. 이는 $\mathbf{C}_t(x,x)$가 외부 지점에도 상당한 가중치를 부여하기 때문이다. $u$-가중으로의 정정은 $\mathbf{C}_t$의 유일한 술어 역할을 제거했다.

$\mathbf{C}_t$는 파생 진단으로 유지된다: 형성 내 지점 쌍과 경계 횡단 지점 쌍을 3자릿수 이상의 차이로 구분하는 능력을 가진다.


\section{부피 제약의 존재론적 지위}
\label{app:cn-volume}

\begin{center}
\fbox{\parbox{0.85\textwidth}{\centering
\textbf{CN6. 부피 제약 $\sum u_i = m$은 구조적 공리이다}\\[4pt]
이것은 계산적 편의가 아니라 유한한 응집 용량에 대한\\
존재론적 약속이다.
}}
\end{center}

\textbf{기술적 필요성.} 부피 제약이 없으면 자명한 해 $u \equiv 0$이 전역 에너지 최소화자이다 (모든 에너지 항이 비음이고 $u = 0$에서 0). 비자명 형성의 존재(T1, T8)에는 부피 제약이 필수적이다.

\textbf{존재론적 자연스러움.} 인지 시스템은 유한한 자원을 가진다. 동시에 무한한 ``응집적 주의''를 기울일 수 없다. 부피 제약 $\sum u_i = m$은 이 유한성의 수학적 표현이다. 부피 분율 $c = m/n$은 ``평균적으로 각 지점에 얼마나 많은 응집적 참여가 허용되는가''를 나타낸다.

\textbf{스피노달 제약.} 위상 전이(T8)는 $c$가 스피노달 범위 $\big(\frac{3-\sqrt{3}}{6}, \frac{3+\sqrt{3}}{6}\big)$ 내에 있어야 발생한다. 이 범위 밖에서는 이중 우물 포텐셜의 2차 도함수 $W''(c) \geq 0$이 되어, 균일 해가 안정하고 형성이 출현하지 않는다.


\section{순서쌍 합 규약}
\label{app:cn-summation}

\begin{center}
\fbox{\parbox{0.85\textwidth}{\centering
\textbf{CN7. 모든 $\sum_{x,y}$는 순서쌍을 범위로 한다}\\[4pt]
이 규약은 위상 전이 정리(T8-Core)와 Hessian 분석에서 핵심적이다.\\
혼동은 인자 2 오류를 발생시킨다.
}}
\end{center}

대칭 핵 $\mathbf{N}(x,y) = \mathbf{N}(y,x)$에서:
\begin{align}
    \text{순서쌍:}\quad &\sum_{x,y} \mathbf{N}(x,y)(u(x) - u(y))^2 = 2\, u^T L\, u \\
    \text{비순서쌍:}\quad &\sum_{\{x,y\}} \mathbf{N}(x,y)(u(x) - u(y))^2 = u^T L\, u
\end{align}

경계/형태학 에너지의 매끄러움 항은 순서쌍 합이므로:
\begin{equation}
    \mathcal{E}_{\mathrm{bd,smooth}} = \alpha \cdot 2u^T L u
\end{equation}

이의 기울기는 $4\alpha L u$ (인자 4, 2가 아님)이다. Hessian 기여도 $4\alpha L$이다. 이 인자들이 위상 전이 임계비에 정확히 진입하므로, 규약 혼동은 임계비를 2배 잘못 계산하게 된다.


\section{자기-참조성의 이중 모드 구조}
\label{app:cn-self-referential}

\begin{center}
\fbox{\parbox{0.85\textwidth}{\centering
\textbf{CN8. 자기-참조성은 이중 모드이다}\\[4pt]
자기-완성(닫힘)과 자기-대비(구별)는 구조적으로 독립인\\
두 가지 자기-참조 양식이다.
}}
\end{center}

SCC의 수학적 독자성은 범용적 비선형성에 있지 않다---모든 비선형 변분 문제는 자기-참조적이다. 독자성은 두 가지 독립적 자기-참조 양식이 동시에 작동하며 서로 다른 구조적 채널을 통해 이론에 진입한다는 데 있다:

\begin{enumerate}
    \item \textbf{자기-완성} (닫힘 $\mathrm{Cl}_t$): 장이 자신의 값을 이용하여 자기-지지를 완성. $\mathcal{E}_{\mathrm{cl}}$과 $\mathsf{Bind}$를 통해 진입.
    \item \textbf{자기-대비} (구별 $\mathbf{D}_t$): 장이 자신의 보(complement) $1{-}u$와의 비대칭을 평가. $\mathcal{E}_{\mathrm{sep}}$와 $\mathsf{Sep}$를 통해 진입.
\end{enumerate}

제3 양식---자기-통합 (공동-소속 $\mathbf{C}_t$)---은 파생 진단으로 이용 가능하나, 현재 어떤 술어나 에너지에도 진입하지 않는다.


\section{다중 형성은 운동학적이다}
\label{app:cn-kinetic}

\begin{center}
\fbox{\parbox{0.85\textwidth}{\centering
\textbf{CN9. 다중 형성은 열역학적이 아니라 운동학적이다}\\[4pt]
연결 그래프에서 전역 최소는 항상 $K^* = 1$이다.\\
$K > 1$은 운동학적 장벽에 의한 준안정 상태로 공존한다.
}}
\end{center}

\textbf{실험적 근거} (exp51): 10가지 그래프 구성(격자, SBM, 바벨, 무작위 기하)에서 $K^* = 1$이 보편적으로 확인되었다. $K > 1$의 공존은 장벽 높이 $\propto \beta^{0.89}$ (exp38, exp55)에 의해 유지되며, 노이즈 $\sigma \leq 0.5$에서 5000번 반복 동안 병합 0회 (exp55).

\textbf{세 기둥}:
\begin{enumerate}
    \item \textbf{핵생성}: 라플라시안 스펙트럼 구조가 초기 형성 분리를 결정
    \item \textbf{준안정성}: 분리 거리 $d > d_{\min}^*$에서 양정치 Hessian
    \item \textbf{조대화}: 노이즈 구동 병합, $K \to (K{-}1)$ 장벽 $\propto \beta^{0.89}$
\end{enumerate}

\textbf{닫힘의 역할} (CN14): 닫힘은 $d_{\min}^*$를 $\sim 30\%$ 감소시킨다 (exp57: $\approx 7 \to \approx 5$ 노드, $\beta = 30$). T7-Enhanced 준안정성의 정량적 실현.
